% BÀI TOÁN HÌNH HỌC - BAI_HINH_28 (TRÍCH ĐỀ 28)
\begin{ex}
Cho đường tròn (O) có đường kính AB và điểm C thuộc đường tròn đó (C khác A, B). Lấy điểm D thuộc dây BC (D khác B, C). Tia AD cắt cung nhỏ BC tại điểm E, tia AC cắt tia BE tại điểm F.
	\begin{enumerate}
		\item[a)] Chứng minh rằng 4 điểm F, C, D, E thuộc một đường tròn.
		\item[b)] Chứng minh rằng $DA \cdot DE = DB \cdot DC$. Gọi I là tâm đường tròn ngoại tiếp tứ giác FCDE, chứng minh rằng IC là tiếp tuyến của đường tròn (O).
	\end{enumerate}
\end{ex}
